\section{Introduction}
\label{sec:intro}

The past three years have produced a series of technical advances in
verifiable AI accountability. Compile-time evidence binding via linear
resource types~\cite{btv2026} ensures that no AI decision can be
issued without a cryptographically bound evidence record. Verifiable
persistence via append-only transparency logs~\cite{btv-t-repo} ensures
that no decision can be delivered without a permanent audit trail.
Privacy-preserving redaction via ZK-SNARKs~\cite{btv-ar2026} ensures
that the right to erasure cannot be weaponised to selectively destroy
evidence of discrimination. Together, these three layers form a
complete accountability stack: no decision without evidence, no
delivery without persistence, no redaction without statistical
justification.

Yet adoption of any component of this stack by production AI systems
remains negligible. This paper investigates why, and whether the
obstacle is economic.

The dominant narrative in AI governance circles is that accountability
tools are too expensive for deployment at scale. If that narrative is
correct, the appropriate policy response is subsidy or mandate.
If it is incorrect --- if transparency infrastructure is, at scale,
cheaper than the regulatory penalties it helps avoid --- then the
appropriate response is information: helping compliance officers and
executives understand the actual cost structure. These two policy
responses are fundamentally different, and choosing between them
requires empirical data.

This paper provides that data. We make four contributions:

\begin{enumerate}
  \item A \textbf{total-cost-of-ownership (TCO) model}
  (Section~\ref{sec:tco}) derived directly from the measured
  benchmarks of the BTV accountability stack, covering infrastructure,
  storage, cryptographic proving, and operations at three deployment
  scales ($10^4$, $10^6$, $10^8$ decisions/year).

  \item An \textbf{empirical penalty baseline}
  (Section~\ref{sec:crossover}) constructed from GDPR enforcement
  data (\euro{}4.5B, 2018--2025~\cite{gdpr-fines-2025}), SEC
  record-keeping fines (>\$2B, 2021--2024~\cite{sec-fines-2024}),
  and projected EU~AI~Act penalties, used to compute the
  \emph{compliance crossover point}: the decision volume above which
  the stack's annualised cost is below expected penalty exposure.

  \item A \textbf{qualitative adoption-barrier analysis}
  (Section~\ref{sec:interviews}) based on $N=12$ semi-structured
  interviews with compliance officers and data protection officers
  (DPOs) at financial institutions and healthcare providers in Brazil
  and the EU, identifying and mapping the three dominant barriers to
  concrete technical mitigations.

  \item A \textbf{compliance-credit mechanism}
  (Section~\ref{sec:credits}) --- a market instrument analogous to
  carbon credits~\cite{ellerman2010} --- whereby operators who deploy
  certified transparency infrastructure receive regulatory safe-harbour
  recognition, converting the cost of accountability into a
  transferable compliance asset and creating an adoption incentive
  without requiring new legislation.
\end{enumerate}

\noindent This paper is the fourth in the BTV programme of research.
Where Papers~1--3 answer ``how can AI accountability be enforced?'',
Paper~4 answers ``under what economic and institutional conditions will
it be adopted?'' The two questions are equally necessary for any
practical deployment.
